\scanpage{119}
\mspage{118}


2) $\unsure{\otimes|\Omega}$-Foncteur $F:C\handwritten{\longrightarrow}C'$ (avec $C,C'$ des trucs-caté\rlap{t}gories).

Les données sont un foncteur \deleted{add} $F$, plus \deleted{d}\unsure{\rlap{e}u|u}\rlap{u}n isomorphisme\deleted{s} de \unsure{bifoncteur|bifoncteurs}
\[
c(X,Y):F(X\otimes Y)\handwritten{\xrightarrow{\sim}}F(X)\otimes F(Y)
\]
\deleted{et isom.}
\[
\deleted{i:F(\underline{1}_C)\unsure{\text{xxx}|\longrightarrow}\underline{1}_{C'}}
\]
satisfaisant les conditions:

a) Compatibilité \inserted{de $c$} avec la donnée de $\handwritten{\phi,\phi'}$

b) Compatibilité \inserted{de $c$} avec la donnée de $\handwritten{\unsure{\mu|\psi},\mu'}$

c) \deleted{Compatibilité de $i$ avec et $c$}\inserted{Existe isom.\ \deleted{fonct} $i\handwritten{\unsure{\rlap{(}:|:}}F(\underline{1})\handwritten{\simeq}\underline{1}'$ compatible avec $c$ et\handwritten{\begin{tikzpicture}[x=.65cm,y=.65cm,baseline=-.2ex]\draw (0,.28) .. controls (.3,.08) and (.77,.24) .. (.93,.05) .. controls (1.21,-.4) and (1.88,-.18) .. (2.26,.1)--(2.25,.7);\draw[line width=.15pt] (-.2,-.5) .. controls (-.35,-.13) and (-.34,.01) .. (-.12,.17);\end{tikzpicture}}} avec les isomorphismes $u\inserted{u'\handwritten{\begin{tikzpicture}[x=.4cm,y=.4cm,baseline=-.2ex]\draw (0,-.25)--(0,.21) .. controls (.4,.48) and (.5,.83) .. (.75,.87) .. controls (.98,.9) and (1.08,1.09) .. (1.1,1.35);\end{tikzpicture}}}$ (ou ce \deleted{qu} qui revient au même, avec \inserted{$c$ et\handwritten{\begin{tikzpicture}[x=.45cm,y=.45cm,baseline=-.2ex]\draw (0,.7) .. controls (.25,.47) and (.72,.67) .. (.99,.35) .. controls (1.08,.16) and (1.14,-.08) .. (1.1,-.23);\draw (1.1,-.23)--(1.55,.88);\end{tikzpicture}}} $v,v'$).

NB Cet $i$ sera alors uniquement déterminé.

d) $F$ est additif.

3) Homomorphisme d'un $\otimes$-foncteur $F$ dans un autre $G$: c'est un homomorphisme de foncteurs $f:F\handwritten{\longrightarrow}G$ satisfaisant l\rlap{e}a\deleted{s} condition

a) compatibilité avec \deleted{$c$, et} $c_F,c_G$

\handwritten{\deleted{\unsure{[a)|(A)} \unsure{pour} \illegible{2}}
\begin{tikzpicture}[x=.65cm,y=.65cm,baseline=-.2ex]
\draw[line width=.2pt] (0,-.25) .. controls (.13,.1) and (.3,.08) .. (.48,.07)--(2.62,.21) .. controls (2.9,.25) and (3.08,.3) .. (3,.48);
\end{tikzpicture}}

\deleted{b)} \deleted{\unsure{(a)|(A)}}\handwritten{(}compatibilité avec $i_F,i_G$ (\deleted{c'}\handwritten{\unsure{ce n'|car}} est \deleted{peut-être}\inserted{\handwritten{\deleted{\unsure{sans doute}} sans doute pas une}} conséquence formelle de a) - à regarder\deleted{)}\handwritten{; \makebox[0pt][l]{\raisebox{1.4ex}{\hspace{.1em}\begin{tikzpicture}[x=1em,y=1em]\fill (0,0) circle (.035);\end{tikzpicture}}}cf plus bas)}

\handwritten{4)} \uline{Rem\unsure{\rlap{\raisebox{1.2ex}{\hspace{.25em}\begin{tikzpicture}[x=1em,y=1em]\draw[line width=1pt] (0,0)--(0,.17);\end{tikzpicture}}}a|a}rque\handwritten{s}}\qquad A dissocier \inserted{\handwritten{ce qui précède\begin{tikzpicture}[x=.55cm,y=.55cm,baseline=-.2ex]\draw (0,-.35) .. controls (-.1,.21) and (.15,.41) .. (.48,.42)--(2.36,.51) .. controls (2.54,.53) and (2.55,.63) .. (2.5,.79);\end{tikzpicture}}} en plusieurs aspects de la donnée d'un bifoncteur\handwritten{$\otimes$}:

a) Donnée d'associativité, soumise à l'axiome du pentagone. Un foncteur ``compatible avec $\otimes$\inserted{associat\deleted{t}if''\handwritten{\begin{tikzpicture}[x=.5cm,y=.5cm,baseline=-.2ex]\draw (0,-.3)--(0,.27) .. controls (.55,.36) and (1.56,.29) .. (2.15,.38) .. controls (2.4,.43) and (2.55,.62) .. (2.55,.82);\end{tikzpicture}}} de $C$ dans $C'$ est un foncteur $+$ un $c$ , compatible avec la donnée de $\handwritten{\phi}$ .

b) Donnée d'unité bilatère $\underline{1},u,v$ (unique à isomorphisme unique près). Un foncteur ``compatible avec \inserted{\handwritten{$\otimes$ à\begin{tikzpicture}[x=.45cm,y=.45cm,baseline=-.2ex]\draw (0,-.25) .. controls (.28,.14) and (.51,.46) .. (.78,.51) .. controls (1.03,.57) and (1.32,.73) .. (1.4,.93);\end{tikzpicture}}} unité\deleted{s} bilatère\deleted{s}'' est un $F$ muni de $c$, tel qu'il existe $i$ compatible avec $c,u,v$ et $c',u',v'$ (lequel $i$ sera alors unique).

\margin{\handwritten{\uline{NB} Le monoïde des endom. de $\underline{1}$ est alors comm.}}

\handwritten{c)} Donnée de commutativité $\handwritten{\mu}$ , satisfaisant la condition de réflexivité.\deleted{(} Un foncteur
